\documentclass[11pt,a4paper]{article}
\usepackage[T1]{fontenc}
\usepackage{lmodern,microtype,amsmath,amssymb,amsthm}
\usepackage[margin=28mm]{geometry}
\newtheorem{proposition}{Proposition}
\newtheorem{corollary}[proposition]{Corollary}
\title{Sharper bounds in count-specific quartic certificates}
\author{Continuation of the count-specific quartic method}
\date{5 November 2026}
\begin{document}
\maketitle
The normal term in the count comparison depends on the same probability
through both its numerator and variance. Bounding these two quantities
separately loses information. The following elementary monotonicity gives
their joint minimum and works on both sides of $p=1/2$.
Throughout, $\Phi$ and $\phi$ are the standard normal distribution function
and density, and $C_{\mathrm E}=(3+\sqrt{10})/(6\sqrt{2\pi})$.
Thresholds $t$ here are in Bernoulli sum coordinates; in a
count cell $t=k+u$, not the standardized threshold of the original sum.

\begin{proposition}[Exact endpoint minimum]
Let $n>0$, $0<l\le h<1$, $0\le t_-\le t_+\le n$ and
$0\le L\le U$. Set
\[
 Z(p,t,\lambda)=\frac{t-np}{\sqrt{np(1-p)+\lambda}},
 \qquad a=t_--nh,\qquad v_h=h(1-h).
\]
Then
\[
 \min_{\substack{p\in[l,h],\ t\in[t_-,t_+]\\\lambda\in[L,U]}}
 \Phi(Z(p,t,\lambda))
 =\Phi\left(\frac{a}{\sqrt{nv_h+\lambda_*}}\right),
 \quad
 \lambda_*=
 \begin{cases} L,&a<0,\\ U,&a\ge0.\end{cases}
 \tag{1}
\]
For a fixed noise $\lambda$, the minimum is at $p=h,t=t_-$.
\end{proposition}
\begin{proof}
Put $D=np(1-p)+\lambda>0$. Differentiation gives
\[
 \frac{\partial Z}{\partial p}
 =-\frac{n\{np+t(1-2p)+2\lambda\}}{2D^{3/2}},
 \qquad \frac{\partial Z}{\partial t}=D^{-1/2}>0.
 \tag{2}
\]
For $p\le1/2$, the bracket is at least $np$ because $t\ge0$.
For $p\ge1/2$, its minimum over $0\le t\le n$ is at $t=n$,
giving the lower bound $n(1-p)$. Thus the first derivative is
strictly negative. This proves the fixed-noise assertion. The remaining
function $a/\sqrt{nv_h+\lambda}$ increases with $\lambda$ when $a<0$
and decreases when $a>0$; it is zero when $a=0$. Monotonicity of $\Phi$
proves (1).
\end{proof}

\begin{corollary}[Sharper whole-cell concavity test]
Under the preceding assumptions, suppose
\[
 t_+-nl\le0,\qquad (t_--nh)^2\le3(nh(1-h)+L).
 \tag{3}
\]
For every fixed $p\in[l,h]$ and $t\in[t_-,t_+]$, the function
$g(\lambda)=\Phi(Z(p,t,\lambda))$ is concave for all $\lambda\ge L$.
\end{corollary}
\begin{proof}
The first condition gives $Z\le0$. Proposition 1 shows that on
$\lambda\ge L$ the smallest value is
$(t_--nh)/\sqrt{nh(1-h)+L}$, whose square is at most three.
Hence $0\le Z^2\le3$ everywhere. With $b=t-np$,
\[
 g''(\lambda)=\frac{b\phi(b/\sqrt D)}{4D^{5/2}}
                    \left(3-\frac{b^2}{D}\right)\le0.
\]
This includes zero noise and the constant case $b=0$; $p\in(0,1)$
keeps $D$ positive.
\end{proof}

\paragraph{Use in an exact certificate.}
For a fixed probability/phase cell, write $t_-=k+u_-$ and
$t_+=k+u_+$. When (3) holds, the quartic upper expectation minus the
normal CDF is convex in the common variance pair $(x,y)$. At a polygon
vertex, set $\lambda=n(x+y)$ and use the directed rational lower bound
for the normal value in (1) with this \emph{same} noise. The bound is
constant on the probability/phase cell for that vertex, so the polynomial
degree and Bernstein acceptance criterion are unchanged. Subtracting the
fixed affine same-law absolute-cube price preserves the same convexity.

When (3) fails, (1) still gives a valid constant Gaussian floor on the
whole noise interval. The concavity claim is then not used. If the
threshold cell is not contained in $[0,n]$, the reader must instead use
the earlier sign-safe variance extrema, or verify a separate derivative
gate. It cannot use (1) without a justification: the bracket in (2) may
be negative for thresholds outside that interval.

The other scalar enclosures also work on either probability half. On
$[l,h]$, use
\begin{align*}
 v_{\min}&=\min\{l(1-l),h(1-h)\},\\
 v_{\max}&=\begin{cases}
 1/4,&l\le1/2\le h,\\
 \max\{l(1-l),h(1-h)\},&\text{otherwise}.
 \end{cases}
\end{align*}
For the acquired mean bound $m\ge m_->0$, put $f(v)=v^2/m_-^2-v$.
Then the noise ceiling satisfies
\[
 \lambda\le n\max\{f(v_{\min}),f(v_{\max})\}.
\]
This quadratic is convex, so its two endpoint values suffice even
without a monotonicity gate. Intersect with the original noise ceiling;
a failed or degenerate intersection is not automatically an exclusion.
For acquired sign gaps $\varepsilon_0,\varepsilon_1>0$, a directed lower
root of $v_{\min}+L/n$ gives the fixed inward ceilings
\[
 r_0=\min(d,h-\varepsilon_0\sqrt{v_{\min}+L/n}_{\,\rm lo}),\quad
 r_1=\min(d,1-l-\varepsilon_1\sqrt{v_{\min}+L/n}_{\,\rm lo}).
\]
They still use the original acquired support and mean hypotheses. This
note proves scalar refinements; acquisition remains a separate input.

\paragraph{Fixed prices retain further acquired moment information.}
In the common-law quartic criterion, put $q=1-p$, $v=pq$,
$x=q\mathbb E U_0^2$, $y=p\mathbb E U_1^2$, $w=v+x+y$ and
$\gamma_i\mathbb E U_i^2=\mathbb E U_i^3$. The exact raw third moments are
\begin{align*}
 R&=\mathbb E|Y-p|^3
   =v(1-2v)+3px+3qy-\gamma_0x+\gamma_1y,\\
 G&=\mathbb E(Y-p)^3
   =v(1-2p)-3px+3qy+\gamma_0x+\gamma_1y.
\end{align*}
These follow by expanding the two conditional cubes. On actual laws,
$R=\rho w^{3/2}$ and $G=\gamma w^{3/2}$, where $\gamma=\mathbb E X^3$.
An acquired interval for each standardized moment and a positive interval
for $w$ give fixed outward bounds $R_-\le R\le R_+$ and
$G_-\le G\le G_+$. For the signed product, retain all four endpoint
products: a negative lower endpoint of $\gamma$ uses an upper width
factor to bound $G$ below.

Let $T$ denote an already justified sufficient CDF target, including
the normal term and any proved lower bound on $C_{\mathrm E}\rho/\sqrt n$.
For fixed rational prices $a,b,c,d\ge0$, replace it by
\[
 T'=T+a(R_--R)+b(R-R_+)+c(G_--G)+d(G-G_+).
 \tag{4}
\]
Every added term is nonpositive on an actual law. Hence proving a
quartic upper expectation strictly below $T'$ suffices. The added
terms are affine in the common $\gamma_i$; at an inward endpoint they
add only linear terms, preserving the previous convex quadratic bound.
After the nine common endpoint choices they are affine in $x,y$,
preserving the variance-polygon convexity. Artificial corners may
violate the acquired moment intervals and receive a higher target;
the inequalities were justified on actual laws before enlargement.
Clearing by $v^2$ gives degree at most eight in $p$, so the original
degree $(9,4)$ suffices for $n=5$. Prices must be fixed on the whole
cell, nonnegative, and shared across all counts and moment choices.

\paragraph{A tangent price for the actual third absolute moment.}
A constant multiple of $R$ can substantially underprice the
normalization at small noise. Convexity gives a more accurate lower
bound. Suppose $0<w\le W$, with $W$ fixed on the cell. The tangent
line to $s^{-3/2}$ at $W$ yields
\[
 w^{-3/2}\ge\frac{5W-3w}{2W^{5/2}},\qquad 5W-3w\ge2W>0.
 \tag{5}
\]
Choose a positive rational $c$ with
$c\le C_{\mathrm E}/(2\sqrt n\,W^{5/2})$ by directed arithmetic.
Since the actual $R$ is nonnegative,
\[
 \frac{C_{\mathrm E}\rho}{\sqrt n}
       \ge cR(5W-3w).
 \tag{6}
\]
The underlying tangent in (5) has first-order contact at $w=W$;
directed rounding of its coefficient adds a fixed lower-bound error.
It can be blended with a fixed lower bound $c_0\rho_-$ using
any fixed rational $\theta\in[0,1]$, where
$0<c_0\le C_{\mathrm E}/\sqrt n$. Both terms are bounds for the
same actual law.

The new target is quadratic in the common variances, so the old
convexity conclusion must be checked again. Write $F$ for the
mixture's quartic upper expectation after the common support-moment
reduction, and fix one of the nine common skew pairs. For fixed $p$,
write
\[
 R=R_0+\alpha_0x+\alpha_1y,\qquad
 R_0=v(1-2v),\quad \alpha_0=3p-\gamma_0,\quad \alpha_1=3q+\gamma_1.
\]
If $H_F$ is the Hessian of $F$ in $(x,y)$,
the Hessian of the quartic loss minus the tangent price is exactly
\[
 H_F+3c\theta
 \begin{pmatrix}2\alpha_0&\alpha_0+\alpha_1\\\alpha_0+\alpha_1&2\alpha_1\end{pmatrix}.
 \tag{7}
\]
Indeed the only added quadratic term in the loss is
$3c\theta(x+y)(\alpha_0x+\alpha_1y)$. For a constant Gaussian floor,
positive semidefiniteness of (7) proves convexity of the complete
loss. Under the concavity test (3), subtracting the actual Gaussian
adds another convex function, so (7) is still sufficient. The
certificate must check this matrix for each common skew choice and
every probability in its cell. Its two diagonal entries and determinant
give an exact scalar test. The added matrix need not itself be
positive semidefinite; omitting this gate would be an invalid proof.

The preceding common-skew reduction still applies before this check:
the new term is affine in each $\gamma_i$ and, at its inward endpoint,
linear in that endpoint. The nonnegative quadratic coefficient from
each quartic is unchanged. Thus no new choice of conditional laws
is introduced, and the variance-polygon vertices suffice once (7)
has been verified.

On a probability/phase cell the cleared tangent target has $p$ degree
at most ten. Consequently the full tensor degree is
$(\max\{n+4,10\},4)$. Multiplying (7) by the positive factor $v^2$
gives polynomial matrix entries; nonnegative Bernstein coefficients
of its diagonal entries and determinant are sufficient for the
uniform matrix gate. Alternatively a fixed-point certificate can
check the rational matrix directly. These are different scopes:
a successful fixed-point check proves no surrounding rectangle
without the corresponding uniform inequalities and cover.
\paragraph{Keeping probability and phase in the Gaussian lower bound.}
Even the exact endpoint minimum (1) can lose a first-order amount across
a wide probability cell. The quartic expectation still varies with $p$,
so this loss need not reflect a small margin in the actual comparison.
A quadratic Taylor lower bound retains this variation at each variance
vertex without changing the common-law convexity argument.

Fix such a vertex and its noise $\lambda$. Write
$F(p,t)=\Phi(Z(p,t,\lambda))$ on a closed rectangle, with center
$(p_0,t_0)$ and half-widths $\Delta_p,\Delta_t>0$.
Suppose rational enclosures give
\[
 F(p_0,t_0)\ge f_0,\qquad
 |F_p(p_0,t_0)-b_p|\le\epsilon_p,\qquad
 |F_t(p_0,t_0)-b_t|\le\epsilon_t,
\]
and $|F_{pp}|\le M_{pp}$, $|F_{pt}|\le M_{pt}$,
$|F_{tt}|\le M_{tt}$ throughout the rectangle.
For any rational $\eta>0$, put
\[
 P=M_{pp}+\eta M_{pt},\qquad T=M_{tt}+M_{pt}/\eta.
\]
Taylor's formula with integral remainder gives, for
$\delta_p=p-p_0$ and $\delta_t=t-t_0$,
\begin{align*}
 F(p,t)\ge{}& f_0+b_p\delta_p+b_t\delta_t
 -\epsilon_p\Delta_p-\epsilon_t\Delta_t\\
 &\hspace{12mm}-\tfrac12(P\delta_p^2+T\delta_t^2).
 \tag{8}
\end{align*}
Indeed the absolute quadratic remainder is at most one half of
$M_{pp}\delta_p^2+2M_{pt}|\delta_p\delta_t|+M_{tt}\delta_t^2$.
Apply $2|ab|\le\eta a^2+b^2/\eta$. The two gradient enclosure
errors cost the displayed fixed budget. The convenient choice
$\eta=\Delta_t/\Delta_p$ balances this mixed error on the rectangle.

Here all required derivative bounds have explicit rational enclosures.
Put $a=t-np$, $D=np(1-p)+\lambda$ and $D'=n(1-2p)$. Directly,
\begin{align*}
 Z_t&=D^{-1/2},&
 Z_p&=-nD^{-1/2}-\frac{aD'}{2D^{3/2}},\\
 Z_{pt}&=-\frac{D'}{2D^{3/2}},&
 Z_{pp}&=\frac{nD'+na}{D^{3/2}}
                    +\frac{3a(D')^2}{4D^{5/2}},\qquad Z_{tt}=0.
\end{align*}
Choose $D_{\min}=nv_{\min}+\lambda$, a positive rational
$s\le\sqrt{D_{\min}}$, and rational bounds $|a|\le A$, $|D'|\le B$.
Valid upper bounds for $|Z|,|Z_p|,|Z_t|,|Z_{pp}|,|Z_{pt}|$ are
\begin{gather*}
 Z_*=A/s,\qquad T_*=1/s,\qquad
 P_*=n/s+AB/(2D_{\min}s),\\
 P_{**}=\frac{n(B+A)}{D_{\min}s}
             +\frac{3AB^2}{4D_{\min}^2s},\qquad
 T_{**}=B/(2D_{\min}s).
\end{gather*}
Since $F_{ij}=\phi(Z)(Z_{ij}-ZZ_iZ_j)$ and $\phi\le2/5$, we may use
\[
 M_{pp}=\tfrac25(Z_*P_*^2+P_{**}),\quad
 M_{pt}=\tfrac25(Z_*P_*T_*+T_{**}),\quad
 M_{tt}=\tfrac25 Z_*T_*^2,
 \tag{9}
\]
or any outward rational upper enclosures of these expressions.
The center gradient is $\phi(Z_0)\nabla Z_0$. In the selected
calculation, $x=a_0^2/(2D_0)\le1$, so the alternating Taylor
polynomials for $e^{-x}$ through orders 24 and 25 enclose this
density after multiplication by directed bounds for $\phi(0)$.
Directed square roots then enclose both derivatives. The existing
exact normal-CDF enclosure supplies $f_0$.

Use (8) only at the fixed variance vertices, after proving convexity
of the loss with the \emph{actual} Gaussian by (3) and (7).
No concavity of the lower polynomial as a function of the variance
is needed or asserted. With $t=k+u$, (8) is a degree-two polynomial
in $p,u$; multiplication by $v^2$ leaves it within the tensor degree
already specified above. Thus it replaces a constant Gaussian floor
without increasing the full certificate's degree.

These refinements establish sufficient certificate criteria. They do
not assert completion of all acquired rows or re-verification of
their original acquisition.
\end{document}
